Um zu zeigen, dass die spezielle lineare Gruppe $\mathrm{SL}_n(\mathbb{K})$ für $n > 1$ nicht kompakt ist, betrachten wir zunächst die Definition der Gruppe. Die Gruppe $\mathrm{SL}_n(\mathbb{K})$ besteht aus allen $n \times n$ Matrizen mit Einträgen aus einem Körper $\mathbb{K}$, deren Determinante gleich 1 ist. 

Eine topologische Gruppe ist kompakt, wenn sie als topologischer Raum kompakt ist, das heißt, jede offene Überdeckung hat eine endliche Teilüberdeckung. Für $\mathrm{SL}_n(\mathbb{K})$ mit $n > 1$ und $\mathbb{K} = \mathbb{R}$ oder $\mathbb{C}$ ist dies nicht der Fall.

Betrachten wir den Fall $\mathbb{K} = \mathbb{R}$. Die Gruppe $\mathrm{SL}_n(\mathbb{R})$ ist eine nicht-kompakte Lie-Gruppe, da sie nicht beschränkt ist. Um dies zu sehen, betrachten wir die Diagonalmatrizen der Form

\[
A_t = \begin{pmatrix}
t & 0 & \cdots & 0 \\
0 & t^{-1} & \cdots & 0 \\
0 & 0 & \ddots & 0 \\
0 & 0 & \cdots & 1
\end{pmatrix}
\]

für $t > 0$. Diese Matrizen liegen in $\mathrm{SL}_n(\mathbb{R})$, da ihre Determinante gleich $t \cdot t^{-1} \cdot 1 \cdots 1 = 1$ ist. Für $t \to \infty$ divergiert die Norm von $A_t$, was zeigt, dass $\mathrm{SL}_n(\mathbb{R})$ nicht beschränkt ist und somit nicht kompakt.

Ein ähnliches Argument gilt für $\mathbb{K} = \mathbb{C}$. Auch hier kann man Matrizen konstruieren, deren Norm unbeschränkt ist, was die Nichtkompaktheit von $\mathrm{SL}_n(\mathbb{C})$ zeigt.

Daher ist $\mathrm{SL}_n(\mathbb{K})$ für $n > 1$ nicht kompakt.